% Report Module 3 – Network Scanning, Attack Classification, and Vulnerability Analysis
% CS 475: Introduction to Computer Security
% Mária Nyolcas
% March 12th 2026

\documentclass[11pt,a4paper]{article}

\usepackage[margin=2.5cm]{geometry}
\usepackage{setspace}
\usepackage{booktabs}
\usepackage{tabularx}
\usepackage{array}
\usepackage{parskip}
\usepackage{titlesec}
\usepackage{hyperref}
\usepackage{enumitem}
\usepackage{microtype}
\usepackage{fancyhdr}
\usepackage{graphicx}
\usepackage{xcolor}

\setstretch{1.15}

\hypersetup{
    colorlinks=true,
    linkcolor=black,
    urlcolor=blue,
    citecolor=black,
}

\titleformat{\section}{\large\bfseries}{}{0em}{}[\titlerule]
\titleformat{\subsection}{\normalsize\bfseries}{}{0em}{}

\pagestyle{fancy}
\fancyhf{}
\lhead{\small CS 475 – Mária Nyolcas}
\rhead{\small Week 3: Network Scanning \& Vulnerability Analysis}
\cfoot{\thepage}
\renewcommand{\headrulewidth}{0.4pt}

\begin{document}
%  TITLE

\begin{center}
  {\Large\bfseries Report Module 3}\\[4pt]
  {\large Network Scanning, Attack Classification, and Vulnerability Analysis}\\[6pt]
  {\normalsize CS 475: Introduction to Computer Security}\\[2pt]
  {\normalsize Mária Nyolcas \quad|\quad Prof: Asen Nikolov Grozdanov}\\[2pt]
  {\normalsize March 2026}
\end{center}

\vspace{4pt}
\noindent\rule{\linewidth}{0.6pt}
\vspace{2pt}

\noindent\textbf{Abstract.}
This report documents a controlled network scanning and vulnerability assessment exercise
conducted entirely within an isolated virtual lab.  A Metasploitable~2 target was
scanned with Nmap using two complementary scan types (SYN and version-detection), three
high-risk attack scenarios were classified against the MITRE ATT\&CK Enterprise framework,
and the top vulnerabilities were identified with OpenVAS, scored with CVSS\,v3.1, and
accompanied by concrete remediation proposals.  Reflections connect the practical findings
to the Week~2 access-control models and consider how defenders can detect the scans performed.


\section{Introduction and Environment Description}

\subsection{Objectives}
The goals of this module are to: (i)~apply host discovery and port-scanning techniques in
a safe lab setting; (ii)~interpret scan output in terms of attack surface; (iii)~classify
realistic attack scenarios using an established taxonomy; and (iv)~perform a basic
vulnerability assessment, prioritise findings by CVSS score, and propose remediations.

\subsection{Lab Topology}
The lab runs on a single Windows host with Hyper-V, hosting two VMs on an isolated
\texttt{host-only} network (\texttt{192.168.56.0/24}) with no routing to the internet
or any production network.

\begin{center}
\small
\begin{tabular}{lllll}
\toprule
Role & OS & IP & Key services & Notes \\
\midrule
Attacker & Kali Linux 2024.3 & 192.168.56.100 & Nmap 7.95, OpenVAS 22 & fully updated \\
Target   & Metasploitable 2  & 192.168.56.101 & many intentional vulns & frozen at release \\
\bottomrule
\end{tabular}
\end{center}

The two machines communicate only over the host-only adapter; no traffic crosses to
external networks during any exercise.

\subsection{Tool Selection}
\textbf{Nmap} was chosen for scanning because its open-source code and exhaustive
documentation make output easy to interpret and reproduce; it also supports both the
half-open SYN mode (useful for stealth comparison) and the \texttt{-sV} version-detection
mode needed for banner grabbing.  \textbf{OpenVAS} (Greenbone Community Edition) was
selected for vulnerability assessment because it is free, actively maintained, and maps
findings directly to CVE identifiers and CVSS scores.  All tool interactions were also
wrapped in custom Python scripts (\texttt{scanner.py}, \texttt{vuln\_assessment.py},
\texttt{attack\_classification.py}) so that parsing and analysis logic is transparent and
reproducible --- following the same philosophy as the Week~2 model simulations.

\section{Scanning Results and Analysis}

\subsection{Host Discovery}
A ping sweep (\texttt{nmap -sn 192.168.56.0/24}) confirmed two live hosts:
\texttt{192.168.56.100} (attacker) and \texttt{192.168.56.101} (target).  ICMP echo
replies were not filtered, which is expected on a host-only adapter with no firewall
policy applied.

\subsection{Port Scan Results}

\textbf{Scan~1 – SYN scan} (\texttt{nmap -sS -p- --open -T4}): sends a TCP SYN, waits
for SYN-ACK (open) or RST (closed), then sends RST to abort the handshake without
completing it.  This is faster and generates less log noise than a full connect.

\textbf{Scan~2 – Version scan} (\texttt{nmap -sV}) on the 12~most interesting ports:
completes the TCP handshake and reads service banners, trading stealth for detail.

\vspace{6pt}
\noindent\small
\begin{tabularx}{\linewidth}{>{\ttfamily}l l l >{\raggedright\arraybackslash}X}
\toprule
\normalfont Port & State & Service & Banner / version (from \texttt{-sV}) \\
\midrule
21/tcp   & open & ftp          & vsftpd 2.3.4 \\
22/tcp   & open & ssh          & OpenSSH 4.7p1 Debian 8ubuntu1 \\
23/tcp   & open & telnet       & Linux telnetd \\
25/tcp   & open & smtp         & Postfix smtpd \\
80/tcp   & open & http         & Apache httpd 2.2.8 (Ubuntu DAV/2) \\
111/tcp  & open & rpcbind      & 2 (RPC \#100000) \\
139/tcp  & open & netbios-ssn  & Samba smbd 3.0.20-Debian \\
445/tcp  & open & microsoft-ds & Samba smbd 3.0.20-Debian \\
3306/tcp & open & mysql        & MySQL 5.0.51a-3ubuntu5 \\
5432/tcp & open & postgresql   & PostgreSQL 8.3.0–8.3.7 \\
6667/tcp & open & irc          & UnrealIRCd 3.2.8.1 \\
8180/tcp & open & http         & Apache Tomcat/Coyote JSP 1.1 \\
\bottomrule
\end{tabularx}
\normalsize

\vspace{6pt}
\noindent\textbf{Analysis.}
The target exposes an unusually large attack surface: 12~ports open on a machine that
should in practice serve only one or two functions.  Several findings stand out:

\begin{itemize}[nosep]
  \item \textbf{Legacy protocols}: Telnet (port~23) and FTP (port~21) transmit
        credentials in cleartext --- a passive network tap captures them without
        any active exploit.
  \item \textbf{Outdated versions}: vsftpd~2.3.4 is known to carry a backdoor
        (CVE-2011-2523); Samba~3.0.20 has a code-injection flaw (CVE-2007-2447);
        Apache~2.2.8 and PHP~5.2 are long past end-of-life.
  \item \textbf{Filtered vs.\ closed}: No ports were returned as \textit{filtered}
        in this environment, meaning no firewall drops packets silently.  In a
        production setting, filtered ports indicate a firewall rule rather than an
        absent service --- useful for mapping defensive perimeter.
  \item \textbf{Detection risk}: The SYN scan is less likely to appear in
        application logs (the handshake is never completed), but it is still
        visible to network IDS tools monitoring for RST floods or SYN without ACK.
        The version scan always completes the handshake and leaves entries in
        every service's connection log.
\end{itemize}

\section{Attack Classification and Vulnerability Findings}

\subsection{MITRE ATT\&CK Mapping}

Three attack scenarios were selected to represent the most realistic threats to the
discovered services.  Each is classified by ATT\&CK tactic (the \emph{why}) and technique
(the \emph{how}), and additionally mapped to STRIDE to show the security property violated.

\vspace{4pt}
\small
\begin{tabularx}{\linewidth}{l >{\raggedright}p{3.2cm} >{\ttfamily}l >{\ttfamily}l >{\raggedright\arraybackslash}X}
\toprule
\# & Service & Tactic (ID) & Technique & STRIDE \\
\midrule
1 & vsftpd 2.3.4 (21) &
  Initial Access\newline TA0001 & T1190 &
  Elevation of Privilege, Information Disclosure \\
2 & Samba 3.0.20 (445) &
  Execution\newline TA0002 & T1059.004 &
  Tampering, Elevation of Privilege \\
3 & OpenSSH / Debian key (22) &
  Credential Access\newline TA0006 & T1110.002 &
  Spoofing, Elevation of Privilege \\
\bottomrule
\end{tabularx}
\normalsize

\vspace{6pt}
\noindent\textbf{Scenario 1 --- vsftpd Backdoor (TA0001 / T1190).}
The compromised 2.3.4 release opens a root shell on port 6200 for any connection whose
username contains `\texttt{:)}'.  This is classified under \emph{Initial Access} because
the attacker's primary objective at this stage is gaining a first foothold.  T1190
(Exploit Public-Facing Application) accurately captures the mechanism: a network-exposed
service is targeted using a publicly known vulnerability.

\noindent\textbf{Scenario 2 --- Samba Command Injection (TA0002 / T1059.004).}
The \texttt{username~map~script} directive is evaluated by \texttt{/bin/sh} without
sanitisation.  An attacker injects shell metacharacters to obtain a reverse shell.
The primary ATT\&CK tactic is \emph{Execution} because the adversary's immediate gain
is code running on the victim; T1059.004 (Unix Shell) names the specific interpreter.

\noindent\textbf{Scenario 3 --- Predictable Debian SSH Keys (TA0006 / T1110.002).}
CVE-2008-0166 reduced the PRNG seed space on Debian to 32,767 values.  Iterating through
a pre-computed key dictionary to authenticate is classified under \emph{Credential Access}
/ Brute Force (T1110), sub-technique T1110.002, because the attacker is effectively
exhausting a small credential space rather than guessing single passwords.

\subsection{Top Vulnerability Findings}

\small
\begin{tabularx}{\linewidth}{>{\ttfamily}l r l >{\raggedright\arraybackslash}X}
\toprule
\normalfont CVE & CVSS\,v3.1 & Severity & Title \\
\midrule
CVE-2011-2523 & 9.8 & CRITICAL & vsftpd 2.3.4 Backdoor RCE \\
CVE-2007-2447 & 9.8 & CRITICAL & Samba Username Map Script RCE \\
CVE-2012-1823 & 9.8 & CRITICAL & PHP CGI Argument Injection RCE \\
CVE-2008-0166 & 7.8 & HIGH     & Debian Predictable OpenSSL PRNG \\
CVE-2004-2761 & 5.3 & MEDIUM   & MD5 Collision in SSL Certificates \\
\bottomrule
\end{tabularx}
\normalsize

\vspace{4pt}
The three CRITICAL findings share the same CVSS base vector
\texttt{AV:N/AC:L/PR:N/UI:N/S:U/C:H/I:H/A:H}: network-reachable, low complexity,
no privileges or user interaction required, and full compromise of confidentiality,
integrity, and availability.  This means any host on the same network (or the internet,
if exposed) can obtain root access with no preconditions --- the highest-possible risk.
CVE-2008-0166 scores 7.8 because it requires local access or prior key enumeration.


\section{Remediation Proposals and Personal Reflection}

\subsection{Remediation}

\noindent\textbf{CVE-2011-2523 --- vsftpd 2.3.4 Backdoor.}
Replace vsftpd immediately with version 3.0.5 or later via the distribution package
manager.  If FTP is not operationally required, remove the service and switch all
file-transfer workflows to SFTP (which is already provided by the OpenSSH daemon on the
same host).  Block port~21/tcp at the firewall as a defence-in-depth measure.

\noindent\textbf{CVE-2007-2447 --- Samba Username Map Script RCE.}
Upgrade Samba to 3.0.25 or any current release.  Remove or blank the
\texttt{username~map~script} smb.conf directive if it is not actively in use.  Restrict
ports 139 and 445 to trusted subnets via firewall rules; SMB should never be exposed to
the internet.  Consider replacing Samba with a dedicated file-share solution if the
full SMB feature set is not needed.

\noindent\textbf{CVE-2012-1823 --- PHP CGI Argument Injection.}
Upgrade PHP to at minimum 5.3.13 / 5.4.3 (or the currently supported 8.x branch).
Switch the web server from CGI mode to PHP-FPM, which does not pass query strings to
the interpreter binary.  As an immediate workaround, add a rewrite rule to strip
query strings beginning with \texttt{-} from PHP script URIs.

\noindent\textbf{Summary table.}

\small
\begin{tabularx}{\linewidth}{>{\ttfamily}l >{\raggedright\arraybackslash}X}
\toprule
\normalfont CVE & Remediation action \\
\midrule
CVE-2011-2523 & Upgrade vsftpd; remove FTP; block port 21 at firewall \\
CVE-2007-2447 & Upgrade Samba; clear username map script; restrict ports 139/445 \\
CVE-2012-1823 & Upgrade PHP; switch to PHP-FPM; add rewrite rule \\
CVE-2008-0166 & Regenerate all SSH keys; audit \texttt{authorized\_keys}; upgrade OpenSSL \\
CVE-2004-2761 & Reissue TLS certificates with SHA-256; disable MD5 cipher suites \\
\bottomrule
\end{tabularx}
\normalsize

\subsection{Personal Reflection}

The most surprising observation was how many critical vulnerabilities are reachable with
zero preconditions.  The vsftpd backdoor requires only a TCP connection and a username
containing two punctuation characters --- an attacker with Metasploit can exploit it in
under thirty seconds.  This made the CVSS scoring feel very real rather than academic:
a 9.8 score genuinely means ``compromise is trivially easy from anywhere on the network.''

Comparing the two scan types reinforced a principle from the Week~2 access-control
discussion: information has value, and generating it has cost.  The SYN scan is faster
and quieter but tells us only which ports are open.  The version scan costs a completed
TCP handshake per port but reveals the exact software version that maps to CVE records
--- the extra information is worth the exposure risk in a controlled assessment.

A limitation I encountered is that OpenVAS produces false positives, particularly for
outdated banner-version checks: the scanner reads the version string in a service banner
and matches it against CVE records without verifying whether the patch has been
back-ported.  I cross-referenced each finding against the NVD entry
(\url{https://nvd.nist.gov}) and the Metasploitable package list to confirm the
vulnerabilities were genuine on this specific image.

Regarding defender detectability: a network IDS such as Snort or Suricata would detect
the SYN scan via a rule matching large numbers of half-open connections to sequential
ports (RST after SYN-ACK pattern).  The version scan is even more visible because each
connection is complete and the banner-grabbing payload is distinctive.  From a defender's
perspective, both scans should trigger an alert within seconds of starting.

Connecting to the broader course arc: the formal models of Week~2 describe what
\emph{should} be enforced; this week reveals what \emph{actually runs} on a system and
how easily misconfigurations and unpatched software undermine those formal guarantees.
A Bell-LaPadula policy is meaningless if a root shell is available via a two-character
backdoor in an FTP daemon.  The coming modules on operating-system security and network
defences will, I expect, examine how to layer technical controls so that even a
misconfigured service does not result in total compromise.


%  REFERENCES
\section*{References}

\begin{enumerate}[label={[\arabic*]}, leftmargin=*, nosep]
  \item MITRE Corporation. (2024). \textit{ATT\&CK Enterprise Matrix v15}. \url{https://attack.mitre.org}
  \item MITRE Corporation. (2024). \textit{CVE Programme}. \url{https://www.cve.org}
  \item FIRST. (2023). \textit{Common Vulnerability Scoring System v3.1 Specification}. \url{https://www.first.org/cvss}
  \item National Institute of Standards and Technology. (2024). \textit{National Vulnerability Database}. \url{https://nvd.nist.gov}
  \item Greenbone Networks. (2024). \textit{Greenbone Community Edition / OpenVAS}. \url{https://www.greenbone.net}
  \item Lyon, G. (2024). \textit{Nmap Network Scanning: The Official Nmap Project Guide}. \url{https://nmap.org/book/}
  \item Rapid7. (2012). \textit{Metasploitable 2 Exploitability Guide}. \url{https://docs.rapid7.com/metasploit/metasploitable-2-exploitability-guide/}
  \item SANS Institute. (2024). \textit{TCP/IP and tcpdump}. \url{https://www.sans.org/posters/}
  \item Microsoft Security. (2002). Shostack, A. \textit{STRIDE Threat Model}. Microsoft Security Engineering.
  \item OWASP Foundation. (2021). \textit{OWASP Top Ten 2021}. \url{https://owasp.org/Top10/}
\end{enumerate}

\end{document}
